\section{Three proof paths and an unconditional black-box tradeoff}
\label{sec:three-paths}

The hybrid admits three logically distinct proof paths.  The first uses only
the published AESP inner-work estimate and the signed LOCSOR tail, and is
closed up to the run-dependent AESP quantity $R$.  The second uses the early
trajectory factor $\Lambda_J$ and becomes optimal under the support-envelope
condition of \cref{sec:support-envelope}.  The third replaces hard confinement
by cumulative KKT-slack charging; its abstract charging inequality is closed,
but two AESP-specific transfer lemmas remain open.

\subsection{A master handoff inequality}

The two tail potentials can be combined without choosing one in advance.

\begin{theorem}[Master handoff inequality]
\label{thm:master-handoff}
Let a Phase-I method return $x_J$ after degree-weighted work $B_J$.  If Phase
II is LOCSOR with $\omega=1$, then
\begin{equation}
\begin{split}
  W_{\rm hybrid}(J)
  \leq B_J+\min\Biggl\{
    &\frac{(1+\alpha)(f(x_J)-f(x^0))}
          {\alpha^2\epsppr^2},\\
    &\frac{1+\alpha}{2\alpha^2\epsppr}M(x_J)
  \Biggr\}.
  \label{eq:master-handoff}
\end{split}
\end{equation}
No sign, confinement, or support-monotonicity assumption is imposed on the
Phase-I trajectory.
\end{theorem}

\begin{proof}
Add the work already spent to the objective-gap bound in
\cref{eq:generic-tail-work} at $\omega=1$ and to the weighted-mass bound in
\cref{eq:mass-tail-general}, then take the smaller certificate.
\end{proof}

The theorem makes the optimal-stopping interpretation precise: AESP pays
$B_J$ to reduce one or both handoff potentials, and LOCSOR pays only for the
remaining certified budget.

\subsection{Path I: an $R$-parameterized theorem without confinement}
\label{subsec:path-I}

Let
\begin{equation}
  c_{\rm A}:=1-\varpi_{\mathrm A},
  \qquad
  \phi_t:=\frac{1+\alpha}{18}c_{\rm A}^t
  \label{eq:path-I-phi}
\end{equation}
be the single-source AESP inner-accuracy schedule.  Recall
$C_t^0=\norm{D^{1/2}\nabla h_t(z_t^{(0)})}_1$, and initialize every stage at
$z_t^{(0)}=y^{(t-1)}$.  Define
\begin{equation}
  R_J:=\max_{1\leq t\leq J}\frac{C_t^0}{\alpha}.
  \label{eq:path-I-R}
\end{equation}
For a single seed, $C_1^0=\alpha$.  Let $B_J^{\rm in}$ denote the inner
adjacency-list work controlled by the AESP theorem; queue construction and
outer sparse-state maintenance are charged separately if they are not already
included.

\begin{proposition}[AESP work to an arbitrary objective handoff]
\label{prop:path-I-burnin}
Fix $\Delta>0$, and let $J\geq1$ be the first stage with $U_J\leq\Delta$,
where $U_t$ is defined in \cref{eq:aesp-outer-bound}.  Then
\begin{equation}
  B_J^{\rm in}
  \leq
  \frac{3000R_J^2\sqrt{\alpha(1-\alpha)}}{\Delta}.
  \label{eq:path-I-burnin}
\end{equation}
\end{proposition}

\begin{proof}
For $\kappa_{\rm A}=1-2\alpha$, the AESP inner contraction parameter is
$\tau=2/3$.
The second term in \cref{eq:aesp-eps-t} gives
\[
  \epsin^{(t)}
  \geq
  \frac{2(1-\alpha)\phi_t}{C_t^0}.
\]
The direct-amortization branch of \cref{eq:aesp-stage-work} therefore yields
\begin{equation}
  W_t^{\rm AESP}
  \leq
  \frac{C_t^0}{\tau\epsin^{(t)}}
  \leq
  \frac{(C_t^0)^2}{2\tau(1-\alpha)\phi_t}
  \leq
  \frac{R_J^2\alpha^2}{2\tau(1-\alpha)\phi_t}.
  \label{eq:path-I-stage}
\end{equation}
Because $J$ is the first crossing,
$U_J>c_{\rm A}\Delta$.  The ratio of the source schedules is
\[
  \frac{U_J}{\phi_J}
  =\frac{3600(1-\alpha)c_{\rm A}}{\alpha},
\]
and hence
\begin{equation}
  \phi_J>\frac{\alpha\Delta}{3600(1-\alpha)}.
  \label{eq:path-I-phiJ}
\end{equation}
Moreover,
\[
  \sum_{t=1}^J\frac1{\phi_t}
  \leq
  \frac1{\varpi_{\mathrm A}\phi_J}.
\]
Sum \cref{eq:path-I-stage}, use \cref{eq:path-I-phiJ}, and substitute
$\tau=2/3$ and
$\varpi_{\mathrm A}=0.9\sqrt{\alpha/(1-\alpha)}$.
\end{proof}

\begin{corollary}[Unconditional $R$-parameterized hybrid bound]
\label{cor:path-I}
Under the conditions of \cref{prop:path-I-burnin},
\begin{equation}
  W_{\rm hybrid}^{\rm in+tail}
  \leq
  \frac{3000R_J^2\sqrt{\alpha(1-\alpha)}}{\Delta}
  +
  \frac{(1+\alpha)\Delta}{\alpha^2\epsppr^2}.
  \label{eq:path-I-tradeoff}
\end{equation}
If an a priori value $\widehat R$ bounds $R_J$ up to the chosen handoff and
\begin{equation}
  \Delta^\star
  :=
  \sqrt{3000}\,
  \frac{\widehat R\,\alpha^{5/4}(1-\alpha)^{1/4}}
       {\sqrt{1+\alpha}}\,\epsppr
  \label{eq:path-I-delta-star}
\end{equation}
lies below the initial AESP certificate, then choosing
$\Delta=\Delta^\star$ gives
\begin{equation}
  W_{\rm hybrid}^{\rm in+tail}
  \leq
  2\sqrt{3000}\,
  \frac{\widehat R(1-\alpha)^{1/4}\sqrt{1+\alpha}}
       {\alpha^{3/4}\epsppr}
  =\bigO\!\left(\frac{\widehat R}{\alpha^{3/4}\epsppr}\right).
  \label{eq:path-I-final}
\end{equation}
\end{corollary}

\begin{proof}
The AESP certificate gives $f(x_J)-f(x^0)\leq\Delta$.  Combine
\cref{prop:path-I-burnin} with the objective branch of
\cref{thm:master-handoff}.  The right-hand side of
\cref{eq:path-I-tradeoff} is $a/\Delta+b\Delta$ and is minimized by
$\sqrt{a/b}$, which is \cref{eq:path-I-delta-star}.
\end{proof}

\begin{remark}[Precise meaning of ``unconditional'']
\Cref{eq:path-I-tradeoff} requires no graph-structural, residual-sign, or
confinement assumption.  The simplified rate in \cref{eq:path-I-final} is
still parameterized by the realized AESP ratio $R_J$.  Since a universal
input-only bound on $R_J$ is not known, the corollary is not yet an absolute
$\bigO(1/(\alpha^{3/4}\epsppr))$ theorem for all instances.  It is nevertheless a
closed instance-wise theorem and improves the direct AESP accuracy dependence
from $\epsppr^{-2}$ to $\epsppr^{-1}$ whenever $R_J$ is controlled.
\end{remark}

\begin{remark}[Tightness of this black-box balancing argument]
Given only a burn-in bound of the form $a/\Delta$ and an objective-tail bound
$b\Delta$, the best switch has value $2\sqrt{ab}$.  Therefore the exponent
$\alpha^{-3/4}$ in \cref{eq:path-I-final} cannot be improved by retuning the
handoff alone; one must strengthen the burn-in estimate, for example through
the early-locality factor in \cref{thm:trajectory}.  This is tightness of the
proof template, not a lower bound for all hybrid algorithms.
\end{remark}

\subsection{Path II: early locality}

Path II is exactly the trajectory theorem in \cref{thm:trajectory}.  It
replaces the inverse-gap burn-in estimate of Path I by
\[
  B_J=\softO\!\left(\frac{\Lambda_J}{\sqrt\alpha}\right).
\]
When $\Lambda_J=\bigO(1/\epsppr)$, the target
$\softO(1/(\sqrt\alpha\epsppr))$ total work follows.  The support-envelope and
no-percolation results in \cref{sec:support-envelope} give one rigorous
sufficient condition; they do not establish the premise on every graph.

\subsection{Path III: cumulative KKT-slack charging}

Hard confinement can be weakened by charging only the volume of coordinates
that become spuriously active.

\begin{lemma}[Abstract KKT-slack charging]
\label{lem:kkt-charge}
Let $\mathcal B\subseteq V$, let every $i\in\mathcal B$ have margin
$\gamma_i\geq\delta>0$, and let $A_k\subseteq\mathcal B$ be the spurious
active set at step $k$.  Suppose activation implies
\[
  \abs{u_k(z_k)_i-u_k(z_k^\star)_i}
  >\gamma_i\sqrt{d_i}
  \qquad(i\in A_k).
\]
Then
\begin{equation}
  \vol(A_k)
  \leq
  \frac{\sqrt{\vol(\mathcal B)}}{\delta}
  \norm{u_k(z_k)-u_k(z_k^\star)}_2.
  \label{eq:kkt-charge}
\end{equation}
If $u_k$ is $L_u$-Lipschitz, then
\begin{equation}
  \sum_k\vol(A_k)
  \leq
  \frac{L_u\sqrt{\vol(\mathcal B)}}{\delta}
  \sum_k\norm{z_k-z_k^\star}_2.
  \label{eq:kkt-path}
\end{equation}
\end{lemma}

\begin{proof}
By Cauchy--Schwarz and $\gamma_i\geq\delta$,
\[
\begin{split}
  \vol(A_k)
  &=\sum_{i\in A_k}
    \frac{\sqrt{d_i}}{\gamma_i}
    \bigl(\gamma_i\sqrt{d_i}\bigr)\\
  &\leq
    \left(\sum_{i\in\mathcal B}\frac{d_i}{\gamma_i^2}\right)^{1/2}
    \left(\sum_{i\in A_k}\gamma_i^2d_i\right)^{1/2}\\
  &\leq
    \frac{\sqrt{\vol(\mathcal B)}}{\delta}
    \norm{u_k(z_k)-u_k(z_k^\star)}_2.
\end{split}
\]
The Lipschitz inequality and summation prove \cref{eq:kkt-path}.
\end{proof}

For the over-regularized RPPR comparison in
\citet{fountoulakis2026complexity}, one obtains a uniform margin
$\delta=\alpha\rho$.  If the relevant accelerated path length is
$\softO(1/\sqrt\alpha)$, \cref{eq:kkt-path} gives the boundary overhead
\[
  \softO\!\left(
    \frac{\sqrt{\vol(\mathcal B)}}{\rho\alpha^{3/2}}
  \right).
\]
To turn this template into an AESP theorem, two transfer lemmas are still
needed: a uniform over-regularized margin for the changing shifted subproblems,
and a summable path-length bound for masked LOCGD or online LOCAPPR inner
iterates.  Thus Path III is a proved accounting mechanism but not yet a
complete graph-uniform AESP bound.
